% ============================================================
% FlexiRefine — Working Paper Draft
% ============================================================
% This is a living draft. See docs/flexirefine_research_log.md
% for raw experiment history, branch log, and result tables.
%
% Last synced: 2026-05-05
% Status: working draft — NOT the Overleaf master copy
%         Do not submit this file directly.
%         Copy sections into Overleaf after advisor review.
%
% Claim status legend used in \todo tags:
%   [CONFIRMED] — supported by 3-seed mean+/-std
%   [SINGLE-RUN] — from one run; directional only
%   [INVALIDATED] — old result superseded; kept for reference
%   [TODO] — needs data or writing
% ============================================================

\documentclass[sigconf, nonacm]{acmart}

\usepackage{amsmath, amssymb}
\usepackage{booktabs}
\usepackage{graphicx}
\usepackage{xspace}
\usepackage{enumitem}
\usepackage{multirow}

\newcommand{\system}{FlexiRefine\xspace}
\newcommand{\quake}{Quake\xspace}
\newcommand{\todo}[1]{\textbf{[TODO: #1]}}
\newcommand{\confirmed}[1]{\textbf{[CONFIRMED: #1]}}
\newcommand{\singlerun}[1]{\textit{[SINGLE-RUN: #1]}}
\newcommand{\superseded}[1]{\textit{[SUPERSEDED: #1]}}

\begin{document}

\title{FlexiRefine: A Control-Plane Framework for Adaptive Maintenance
       in Dynamic Vector Indexing}

\author{Adiyan Islam}
\affiliation{%
  \institution{University of Washington}
  \city{Seattle}
  \country{USA}
}
\author{Dongfang Zhao}
\affiliation{%
  \institution{University of Washington}
  \city{Seattle}
  \country{USA}
}

% ============================================================
\begin{abstract}
% ============================================================
% [CONFIRMED] Main claim updated from old 22% to 36-46% across 3 seeds.
% [CONFIRMED] All selective configs preserve mean recall >= 0.90.
% [CONFIRMED] Total runtime within 4% of full-refinement baseline.

Dynamic vector search systems must maintain high recall and low query
latency under continuous inserts, deletes, and shifting access patterns.
Existing indexes expose only limited control over the tradeoff between
recall, search latency, update cost, and maintenance overhead.
Partitioned indexes are update-friendly but suffer recall loss if
maintenance is disabled; enabling full maintenance can dominate
runtime even under moderate update rates.

We present \system, a control-plane framework that treats refinement---a
key post-split maintenance operation---as a tunable resource allocation
problem. Rather than refining all candidate partitions after a split,
\system ranks candidates by a scoring function that jointly estimates
workload importance (via query hit counts) and structural staleness (via
mutation counts since last refinement). This exposes a controllable
tradeoff between maintenance cost, search quality, and total runtime.

We implement a family of refinement policies including split-threshold
tuning, global and per-split candidate budgeting, hit-based selection,
raw and density-normalized staleness scoring, and a general multiplicative
score family parameterized by importance and staleness exponents.
We evaluate on a dynamic SIFT1M workload with inserts, deletes, and
queries, reporting mean\,$\pm$\,std across three independent workload
seeds.

\textbf{Score-based selective refinement reduces maintenance by
36--46\% while preserving mean recall above 0.90 across three
independent workloads, with total runtime within 4\% of the
full-refinement baseline.}
Naive filtering policies (size-only, distance-only, per-split coverage)
reduce maintenance but fail to preserve the recall target, confirming
that a joint workload-importance--staleness signal is necessary.
\end{abstract}

\maketitle

% ============================================================
\section{Introduction}
\label{sec:intro}
% ============================================================

Vector search is a core primitive in retrieval-augmented generation,
recommendation systems, semantic search, and embedding-based
personalization. These applications store large collections of
high-dimensional vectors and answer approximate nearest-neighbor queries
under strict latency and recall requirements. In production, the indexed
data rarely remains static: new embeddings arrive continuously, stale
vectors are deleted, and query distributions shift with user interests.

Existing vector indexes only partially address this dynamic setting.
Graph-based indexes such as HNSW~\cite{malkov2018efficient} and
DiskANN~\cite{jayaram2019diskann} achieve strong search performance in
read-heavy settings but incur high update cost as graph quality
degrades. Partitioned indexes such as Faiss-IVF~\cite{johnson2019billion}
and ScaNN~\cite{guo2020accelerating} are more update-friendly because
inserts and deletes are local operations, but they suffer from partition
imbalance and recall loss under dynamic workloads if maintenance is
neglected. Recent adaptive systems, including \quake~\cite{todo-quake},
show that cost-model-driven maintenance can significantly improve dynamic
partitioned vector search. However, even these systems expose only a
limited set of operating points: maintenance is either fully enabled or
disabled, with few controllable knobs between these extremes.

The limitation is not only algorithmic---it is also a control problem.
Different users have different tradeoff preferences. A
retrieval-augmented generation system may require recall near 0.90 or
higher and accept higher maintenance overhead. A high-ingestion pipeline
may prefer aggressive maintenance reduction even if recall drops
slightly. A latency-sensitive recommendation service cares most about
stable query latency. A single fixed maintenance policy cannot serve all
of these needs.

This paper argues that dynamic vector indexing systems should expose a
flexible control surface over maintenance decisions. We focus on
\emph{refinement}, the operation that repairs local partition quality
after splits and updates by running K-means on neighboring partitions.
Refinement is important because it reduces overlap and improves vector
assignments, but it is also expensive: in our experiments, disabling
refinement cuts maintenance time by roughly 14$\times$ while dropping
recall by approximately 3 points. This makes refinement a natural
control point---costly enough to matter, important enough that removing
it blindly is unsafe.

\system is a control-plane framework for selective refinement. The key
idea is that refinement is most valuable when applied to partitions that
are both \emph{important to current queries} (approximated by recent
query hit counts) and \emph{structurally stale} (approximated by
mutations since last refinement). Given a maintenance budget, the system
ranks candidate partitions by a score combining these two signals and
refines only the top-ranked subset.

This paper makes the following contributions:

\begin{itemize}[leftmargin=*]
  \item \textbf{Control-plane framing.}
    We frame dynamic vector index maintenance as a controllable
    tradeoff problem, enabling users to select operating points based
    on application priorities rather than accepting a fixed policy.

  \item \textbf{System design and implementation.}
    We design and implement \system as a modular refinement-control
    layer integrated with a dynamic partitioned vector index. The layer
    includes candidate generation, per-partition metadata tracking,
    scoring, budget allocation, and refinement execution. We expose all
    policy knobs through configuration files for reproducibility.

  \item \textbf{Score-family theory.}
    We formulate a general multiplicative score family
    $S_{\beta,\gamma}(p) = (h_p + 1)^\beta D(p)^\gamma$
    that unifies hit-only, mutation-only, raw staleness, and
    density-normalized staleness policies as special cases. We prove
    budget monotonicity and cost bounding properties.

  \item \textbf{Multiseed empirical evaluation.}
    We evaluate a wide range of policies and report mean\,$\pm$\,std
    across three independent workload seeds to account for K-means
    non-determinism. Score-based selective refinement reduces
    maintenance by 36--46\% while preserving mean recall above 0.90
    with total runtime within 4\% of the baseline.

  \item \textbf{Negative-result analysis.}
    We show that naive policies---size-only, distance-only,
    per-split coverage budgeting, and hits-alone---each fail to
    preserve the recall target while reducing maintenance. These
    negative results motivate the joint importance--staleness signal.

  \item \textbf{Reproducibility infrastructure.}
    All experiment configurations are YAML files included in the
    repository. The multi-seed runner produces raw per-seed CSVs and
    aggregated mean\,$\pm$\,std tables automatically.
    \todo{Insert final repository URL when ready.}
\end{itemize}

The rest of the paper proceeds as follows.
Section~\ref{sec:motivation} motivates user-facing tradeoffs.
Section~\ref{sec:system} presents the \system architecture and
methodology.
Section~\ref{sec:theory} gives the theoretical formulation.
Section~\ref{sec:implementation} describes implementation details.
Section~\ref{sec:evaluation} evaluates the framework.
Section~\ref{sec:related} discusses related work.
Section~\ref{sec:discussion} discusses limitations and variance.
Section~\ref{sec:conclusion} concludes.
Section~\ref{sec:future} identifies future work.

% ============================================================
\section{Motivation and Use Cases}
\label{sec:motivation}
% ============================================================

Dynamic vector search is inherently multi-objective. A system must
balance at least four quantities: search latency $T_s$, maintenance
time $T_m$, update time $T_u$, and recall $R$.

A useful application-level objective is:
\begin{equation}
  J = \lambda_s T_s + \lambda_m T_m + \lambda_u T_u - \lambda_r R
\end{equation}
where the $\lambda$ coefficients encode user priorities. This objective
is not meant to imply that every deployment explicitly sets these
weights. It captures the reality that different systems value the same
metrics differently.

\subsection{Recall-Critical Users}
A search or RAG application may require recall near 0.90 or higher.
For this user, $\lambda_r$ is large. The system should choose
conservative policies such as full refinement or high-budget
score-based refinement. Maintenance cost is acceptable if it preserves
accuracy.

\subsection{Maintenance-Constrained Users}
A high-ingestion workload with limited background resources prefers
aggressive selective refinement. Recall around 0.88--0.90 may be
acceptable if maintenance cost drops substantially.

\subsection{Search-Sensitive Users}
A latency-sensitive recommender cares most about stable query latency.
The system should avoid under-refining hot query regions. Hit-based or
high-budget score-based policies are natural choices.

\subsection{Balanced Users}
Many deployments want recall near target, lower maintenance, and no
increase in total runtime. This is the regime where score-based
selective refinement is most useful.

\subsection{Design Goal}
\system makes these choices explicit through four control dimensions:
(1) how many partitions should be refined,
(2) which signals should rank candidates,
(3) how aggressively to prioritize recall versus maintenance savings,
and (4) how the refinement budget should be distributed across split
neighborhoods.

% ============================================================
\section{\system System Design}
\label{sec:system}
% ============================================================

\subsection{System Model}

We consider a dynamic partitioned vector index. The index stores vectors
in partitions, each represented by a centroid. Queries search a subset
of partitions and return approximate nearest neighbors. Updates insert
or delete vectors. Maintenance periodically modifies the index structure
through splits, merges, and refinements.

Let $P$ be the set of partitions. Each partition $p \in P$ has:
\begin{itemize}[leftmargin=*]
  \item size $|p|$,
  \item centroid $c_p$,
  \item recent query-hit count $h_p$,
  \item mutation count $m_p$ (since last refinement),
  \item minimum distance to recent split centroids $d_p$.
\end{itemize}

\system operates during maintenance. When the underlying index decides
that one or more partitions should be split, \system controls the
refinement stage that follows those splits.

\subsection{Architecture Overview}

\system consists of five modules:

\begin{enumerate}[leftmargin=*]
  \item \textbf{Candidate Generator.}
    Given newly split partition centroids, retrieve the $r_f$ nearest
    partition centroids from the parent index for each split centroid.
    Deduplicate across split centroids by keeping the minimum distance
    per candidate partition.

  \item \textbf{Metadata Tracker.}
    Maintain per-partition statistics: size, query hits, mutation count,
    and distance to split centroids. Update after every insert, delete,
    and refinement operation.

  \item \textbf{Scoring Engine.}
    Compute a refinement score for each candidate partition using a
    selectable function of hit count and staleness.

  \item \textbf{Budget Allocator.}
    Select a subset of candidates under a global top-$K$ budget,
    per-split budget, score threshold, or hybrid policy.

  \item \textbf{Refinement Executor.}
    Apply local K-means refinement to selected candidates.
    Reset mutation counters and update parent-index centroids.
\end{enumerate}

\begin{figure}[h]
  \centering
  \fbox{\parbox{0.9\linewidth}{%
    \centering
    \vspace{0.4in}
    \textbf{Placeholder: architecture diagram.}\\
    Query/update stream $\rightarrow$ Metadata Tracker
    $\rightarrow$ maintenance trigger
    $\rightarrow$ Candidate Generator
    $\rightarrow$ Scoring Engine
    $\rightarrow$ Budget Allocator
    $\rightarrow$ Refinement Executor.
    \vspace{0.4in}
  }}
  \caption{High-level \system architecture. The refinement-control layer
    sits between split decisions and refinement execution, leaving the
    base index unchanged.}
  \label{fig:architecture}
\end{figure}

\subsection{Problem Formulation}

After a split, the candidate generator produces a set $C$. A
refinement policy $\pi$ selects a subset:
\begin{equation}
  \pi(C, B) \subseteq C, \quad |\pi(C,B)| \leq B
\end{equation}
The ideal policy minimizes the multi-objective:
\begin{equation}
  \pi^* = \arg\min_{\pi(C,B)}
    \bigl[\lambda_s T_s(\pi) + \lambda_m T_m(\pi) + \lambda_u T_u(\pi)
          - \lambda_r R(\pi)\bigr]
\end{equation}
This is not solvable online because the effect of refining a candidate
on future recall and search time is not known. \system approximates
utility through observable metadata.

\subsection{Candidate Generation}

For each split centroid $s \in S$, retrieve the $r_f$ nearest
partition centroids:
\begin{equation}
  C_s = \mathrm{Nearest}_{r_f}(s,\, P)
\end{equation}
The global candidate set is $C = \bigcup_{s \in S} C_s$, deduplicated
by minimum distance: $d_p = \min_{s \in S} \|c_p - c_s\|$.

\textit{Intuition.} The radius $r_f$ controls the spatial scope of
repair. Distance alone is not sufficient for final selection---it
determines which partitions are geometrically affected, not which
are important to queries or structurally degraded.

\subsection{Metadata Tracking}

\textbf{Query hits.}
For a recent query window $W$ of size $|W|$ queries:
\begin{equation}
  h_p = \sum_{q \in W} \mathbf{1}[p \in \mathrm{scanned}(q)]
\end{equation}
Maintained in a circular sliding-window buffer. All queries in the
window are weighted uniformly. Reuses query-hit data already gathered
by the adaptive partition scanning mechanism; no additional profiling
overhead.

\textbf{Mutation count.}
Each insert or delete affecting partition $p$ increments the counter:
\begin{equation}
  m_p \leftarrow m_p + \Delta_p
\end{equation}
After partition $p$ is refined, the counter is reset:
\begin{equation}
  m_p \leftarrow 0 \quad \text{(after refinement)}
\end{equation}
Newly split child partitions are initialized with
$m_p = \max(1, |p|)$ to ensure they enter the scoring system as
non-trivially stale.

\textbf{Mutation density.}
The density-normalized staleness signal divides the raw mutation count
by current partition size:
\begin{equation}
  \tilde{m}_p = \frac{m_p}{|p|}
  \label{eq:density}
\end{equation}
This estimates the mutation \emph{rate} per resident vector, which
better approximates centroid displacement than absolute mutation count:
a large partition with $m_p = 2000$ and $|p| = 3000$ has lower
centroid drift than a small partition with the same $m_p$ and $|p| = 500$.

\subsection{Refinement Control Policies}

\subsubsection{Split-Threshold Tuning}
A partition $p$ is split when the predicted cost delta is negative:
\begin{equation}
  \Delta C_{\mathrm{split}} < -\tau_s
\end{equation}
Higher $\tau_s$ reduces splits (and therefore total refinement
opportunities). Split-threshold tuning alone reduces maintenance
substantially but cannot recover the recall target without a strong
candidate-ranking signal (see~\S\ref{sec:eval_q2}).

\subsubsection{Refinement Radius}
The radius $r_f$ controls candidate set size:
\begin{equation}
  |C| \leq r_f |S|
\end{equation}
It is a geometric candidate-generation parameter, not a complete
refinement policy. Increasing $r_f$ broadens coverage but also
increases maintenance cost.

\subsubsection{Refinement Iterations}
Let $i_f$ be the number of K-means iterations per refinement call.
Refinement cost scales as $O(i_f \cdot |C'| \cdot \bar{s} \cdot d)$
where $C'$ is the selected subset, $\bar{s}$ is average selected
partition size, and $d$ is vector dimensionality. More iterations
repair structure more strongly but exhibit diminishing returns.

\subsubsection{Size-Based Filtering}
\begin{equation}
  S_{\mathrm{size}}(p) = \mathbf{1}[|p| \geq \tau_{\mathrm{size}}]
\end{equation}
\textit{Expected benefit.} Large partitions cost more to scan.
\textit{Weakness.} Size does not measure query importance or staleness;
this policy refines large stable partitions while ignoring small
stale ones. Empirically shown to be a weak signal (see~\S\ref{sec:eval_q8}).

\subsubsection{Distance-Based Filtering}
\begin{equation}
  p \text{ is selected if } d_p \leq \tau_d
\end{equation}
\textit{Expected benefit.} Nearby partitions are more likely to be
affected by local geometry changes.
\textit{Weakness.} Distance does not measure query importance.
Partitions near a split may be geometrically affected but rarely queried.

\subsubsection{Global Top-K Selection}
\begin{equation}
  \pi(C, K) = \mathrm{TopK}_{p \in C}\; S(p)
\end{equation}
Controls the global refinement budget directly. Required for the
score-based policies to be effective.

\subsubsection{Per-Split Budgeting}
\begin{equation}
  \pi(C, b) = \bigcup_{s \in S} \mathrm{Top\text{-}}b(C_s)
\end{equation}
Ensures every split receives some local coverage by allocating $b$
candidates per split centroid before the global union.
\textit{Weakness (negative result).}
Per-split forcing breaks global value ranking. When split centroids
are spatially close, their neighborhoods overlap; capping per centroid
blocks globally high-value candidates that do not rank first from
every nearby centroid. Empirically, per-split budgeting does not
preserve recall above 0.90 under any tested configuration after the
correct implementation of the per-centroid cap
(see~\S\ref{sec:eval_q5}).

\subsubsection{Hit-Based Refinement}
\begin{equation}
  S_{\mathrm{hit}}(p) = h_p
\end{equation}
\textit{Expected benefit.} Prioritizes query-important partitions.
\textit{Weakness.} Does not account for structural degradation.
A frequently queried but stable partition may consume budget at the
expense of a stale but less-queried one.

\subsubsection{Hits-by-Mutations (Raw Staleness)}
\begin{equation}
  S_{\mathrm{hm}}(p) = (h_p + 1)\, m_p
  \label{eq:raw_score}
\end{equation}
Selects partitions that are both important (high $h_p$) and stale
(high $m_p$). The $+1$ smoothing ensures partition with zero hits
but high mutations still receives a non-trivial score.

\subsubsection{Hits-by-Density (Density-Normalized Staleness)}
\begin{equation}
  S_{\mathrm{density}}(p) = (h_p + 1)\, \frac{m_p}{|p|}
  \label{eq:density_score}
\end{equation}
Replaces raw mutation count with mutation density
(Eq.~\ref{eq:density}). Empirically, this configuration exhibits
lower maintenance variance across seeds than raw staleness scoring at
comparable recall levels (see~\S\ref{sec:eval_q6}).

% ============================================================
\section{Theoretical Formulation}
\label{sec:theory}
% ============================================================

\subsection{Refinement Utility}

Let the expected benefit of refining partition $p$ be:
\begin{equation}
  B(p) = \Delta R(p) - \alpha \Delta T_s(p)
\end{equation}
where $\Delta R(p)$ is the expected recall improvement and
$\Delta T_s(p)$ is the expected search-time reduction.
The net utility is:
\begin{equation}
  U(p) = B(p) - \beta K(p)
\end{equation}
where $K(p) = T_{\mathrm{refine}}(p)$ is the refinement cost.
Since $B(p)$ is not observable online, \system approximates it with
measurable metadata.

\subsection{Importance}

Partition importance estimates future query dependence:
\begin{equation}
  I(p) = f_h(h_p)
\end{equation}
Examples: $f_h = h_p$ (linear), $f_h = \log(1+h_p)$ (diminishing
returns), $f_h = \sigma(\alpha h_p)$ (sigmoid saturation). The log
and sigmoid variants model the intuition that extremely hot partitions
should not dominate the full budget.

\subsection{Staleness}

Partition staleness estimates structural degradation:
\begin{equation}
  D(p) = f_m(m_p)
\end{equation}
Staleness transforms:
\begin{equation}
  f_m(m_p) = m_p \quad \text{(raw)}
\end{equation}
\begin{equation}
  f_m(m_p) = \frac{m_p}{|p|} \quad \text{(density-normalized)}
\end{equation}
\begin{equation}
  f_m(m_p) = \sqrt{m_p} \quad \text{(sqrt)}
\end{equation}
\begin{equation}
  f_m(m_p) = \log(1+m_p) \quad \text{(log)}
\end{equation}
These capture different assumptions about how structural degradation
accumulates as a function of mutation count and partition size.

\subsection{General Score Family}

The general multiplicative family is:
\begin{equation}
  S_{\beta,\gamma}(p) = (h_p + \epsilon)^\beta\, D(p)^\gamma
  \label{eq:general}
\end{equation}
where $\beta$ controls importance sensitivity, $\gamma$ controls
staleness sensitivity, and $\epsilon$ smooths sparse hit counts.
$D(p)$ can be any staleness transform (raw, density, sqrt, log).

Special cases:
\begin{center}
\begin{tabular}{lll}
  \toprule
  $\beta$ & $\gamma$ & $D(p)$ \\ \midrule
  1 & 0 & any & Hit-only \\
  0 & 1 & $m_p$ & Mutation-only \\
  1 & 1 & $m_p$ & Hits-by-mutations (raw) \\
  1 & 1 & $m_p/|p|$ & Density-normalized (Density70-Linear) \\
  0.5 & 1 & $m_p/|p|$ & SqrtHits + density \\
  1 & 0.5 & $m_p/|p|$ & SqrtMut + density \\
  0.5 & 0.5 & $m_p/|p|$ & SqrtBoth + density \\
  \bottomrule
\end{tabular}
\end{center}

An activation-style variant:
\begin{equation}
  S(p) = \sigma(\alpha h_p)\,\log(1+m_p)
\end{equation}
An additive feature model:
\begin{equation}
  S(p) = w_h\phi(h_p) + w_m\psi(m_p) + w_d\chi(d_p) + w_s\eta(|p|)
\end{equation}
\todo{Activation-style and additive variants not yet evaluated experimentally.}

\subsection{Property 1: Budget Monotonicity}

For fixed candidate set $C$ and score function $S$, if $B_1 \leq B_2$:
\begin{equation}
  \mathrm{TopK}(C, B_1) \subseteq \mathrm{TopK}(C, B_2)
\end{equation}
Budget provides a monotone control knob. Increasing $B$ cannot
decrease the number of refined candidates.

\subsection{Property 2: Refinement Cost Upper Bound}

With at most $B$ selected candidates of maximum size $s_{\max}$,
dimension $d$, and $i_f$ refinement iterations:
\begin{equation}
  T_{\mathrm{refine}} = O(B \cdot i_f \cdot s_{\max} \cdot d)
\end{equation}
The user can cap worst-case refinement work by controlling $B$ and
$i_f$ independently.

\subsection{Property 3: Score-Based Selection Approximates Refinement Return}

Suppose the expected benefit of refining $p$ increases monotonically
with both $I(p)$ and $D(p)$. If $S(p) = I(p)\,D(p)$, then high-score
partitions are exactly those with large joint importance and staleness.
The product suppresses candidates that are hot but stable (large $I$,
small $D$) or stale but irrelevant (large $D$, small $I$).

% ============================================================
\section{System Implementation}
\label{sec:implementation}
% ============================================================

\subsection{Code Structure}

\system is implemented as a set of modifications to the maintenance
path of a dynamic partitioned vector index. The implementation
modifies four components:

\begin{itemize}[leftmargin=*]
  \item \textbf{Maintenance policy} (\texttt{maintenance\_policies.cpp}):
    candidate ranking, per-split cap, scoring, and selective refinement.
  \item \textbf{Partition manager} (\texttt{partition\_manager.cpp}):
    per-partition mutation counter, incremented on add/remove,
    reset to zero after refinement.
  \item \textbf{Struct definition} (\texttt{common.h}):
    \texttt{MaintenancePolicyParams} with all tunable fields.
  \item \textbf{Python bindings} (\texttt{wrap.cpp}):
    exposes all new parameters to Python and YAML configuration.
\end{itemize}

\subsection{New Parameters}

\system adds the following maintenance parameters
(all backward-compatible; existing configs unaffected by default values):

\begin{itemize}[leftmargin=*]
  \item \texttt{refinement\_max\_candidates} (int, default $-1$):
    global cap on candidate partitions after deduplication.
  \item \texttt{refinement\_candidates\_per\_split} (int, default $-1$):
    per-split centroid cap applied before global union.
  \item \texttt{refinement\_size\_threshold} (int, default $-1$):
    minimum partition size to be eligible for refinement.
  \item \texttt{refinement\_distance\_threshold} (float, default $-1$):
    maximum centroid distance to any split centroid to be eligible.
  \item \texttt{refinement\_top\_k\_hits} (int, default $-1$):
    keep top candidates by hit count.
  \item \texttt{refinement\_top\_k\_score} (int, default $-1$):
    keep top candidates by combined score.
  \item \texttt{refinement\_normalize\_mutations} (bool, default \texttt{false}):
    if true, staleness $= m_p / |p|$ instead of $m_p$.
  \item \texttt{refinement\_score\_beta} (double, default 1.0):
    exponent on $(h_p + 1)$ in the score formula.
  \item \texttt{refinement\_score\_gamma} (double, default 1.0):
    exponent on staleness in the score formula.
\end{itemize}

With all parameters at defaults, score computation is:
\begin{equation}
  \mathrm{score} = (h_p + 1)^{1.0} \times m_p^{1.0}
    = (h_p + 1) \times m_p
\end{equation}
which is numerically identical to the original formula.

\subsection{Mutation Tracking}

For each partition $p$, a lightweight integer counter
$m_p$ is maintained:
\begin{align}
  m_p &\leftarrow m_p + \Delta_p \quad \text{(on add or remove)} \\
  m_p &\leftarrow 0 \quad \text{(after } \texttt{refine\_partitions}\text{)}
\end{align}
Newly added partitions (children of splits) are seeded with
$m_p = \max(1, |p|)$ so they enter scoring as non-trivially stale.

\subsection{Hit Tracking}

The system reuses query-hit data already gathered for the adaptive
partition scanning mechanism. For each query, the partitions scanned
are recorded. Over window $W$:
\begin{equation}
  h_p = \sum_{q \in W} \mathbf{1}[p \in \mathrm{scanned}(q)]
\end{equation}
Stored in a circular buffer of size $|W|$. Queries are currently
weighted uniformly; recency-weighted variants are possible future work.

\subsection{Selective Refinement Pipeline}

\begin{enumerate}[leftmargin=*]
  \item Split partitions produced by the base maintenance policy.
  \item For each split centroid, retrieve the $r_f$ nearest partition
    centroids from the parent index.
  \item Apply per-split cap (if \texttt{refinement\_candidates\_per\_split} $> 0$):
    sort each centroid's neighbor list by distance, truncate to cap.
  \item Union across centroids into a global min-distance map.
  \item Apply optional distance filter and size filter.
  \item Compute score for each candidate using Eq.~\ref{eq:general}.
  \item Apply top-$K$ budget: keep highest-scoring candidates.
  \item Run K-means refinement on selected candidates.
  \item Reset mutation counters for refined partitions.
\end{enumerate}

\subsection{Engineering Considerations}

\begin{itemize}[leftmargin=*]
  \item Candidate scoring is $O(|C|)$ in the local refinement
    neighborhood, not over all partitions.
  \item Hit counts reuse existing maintenance metadata; no new query
    interception is needed.
  \item Mutation counters are lightweight 64-bit integer fields per partition.
  \item All new parameters are exposed through YAML configuration files
    with backward-compatible defaults.
\end{itemize}

% ============================================================
\section{Evaluation}
\label{sec:evaluation}
% ============================================================

% DZ revision 2026-05-05:
% Evaluation is organized around research questions, not by baseline.
% Each subsection answers one question with a result-style title.
% Research questions are listed explicitly here for orientation.

\noindent\textbf{Research questions.}
Our evaluation addresses the following questions:
\begin{enumerate}[leftmargin=*,label=\textbf{RQ\arabic*.}]
  \item \textit{Does FlexiRefine reduce maintenance while preserving recall?}
        Compares Quake full refinement against \system selective-refinement policies.
  \item \textit{Are score-based policies better than naive selective policies?}
        Evaluates no-refine, split-threshold, top-K, per-split, hit-only,
        and naive-budget baselines.
  \item \textit{Is dynamic maintenance necessary under insert/delete workloads?}
        Compares \system against Faiss-IVF (no dynamic maintenance) and
        LIRE-style (dynamic splits without K-means refinement).
  \item \textit{Are the results robust across workload seeds?}
        Reports mean\,$\pm$\,std across three independent workloads.
  \item \textit{Which operating point should a user choose?}
        Maps user priorities to specific \system configurations.
\end{enumerate}

\subsection{Experimental Setup}
\label{sec:eval_setup}

\textbf{Platform.}
\todo{Add CloudLab machine specification.}

\textbf{Dataset.}
SIFT1M: 1 million 128-dimensional floating-point vectors, Euclidean distance.

\textbf{Workload.}
Dynamic workload with 30\% inserts, 20\% deletes, and 50\% queries,
1000 operations, initial size 100,000 vectors, skewed cluster insert
distribution, uniform query distribution. Target recall 0.90.

\textbf{Seeds.}
Primary confirmation results report mean\,$\pm$\,std across seeds
$\{9299, 42, 12345\}$. Exploratory tables use seed 9299 only and are
labeled \textit{[single-run]}.

\textbf{Baselines.}
\begin{itemize}[leftmargin=*]
  \item \textbf{Quake} --- full cost-model maintenance with complete K-means
    refinement. This is the system \system is built on and the primary
    full-refinement reference point.
  \item \textbf{Faiss-IVF} (nprobe=20) --- static IVF index trained once at
    build time; no dynamic maintenance. Centroids never updated; recall
    degrades as the workload evolves. Included as a no-maintenance
    degradation baseline. nprobe=20 was selected by a tuning sweep to
    achieve mean recall $\approx$ 0.90 at workload start.
  \item \textbf{LIRE-style} --- Quake backend with size-based splits
    (\texttt{max\_partition\_size}=2000) and no K-means refinement
    (\texttt{refinement\_iterations}=0, \texttt{refinement\_radius}=0).
    Approximates SPFresh/LIRE design: dynamic partitioning without
    centroid-quality repair.
    \todo{Replace with confirmed 3-seed LIRE results when available.}
  \item \textit{HNSW, ScaNN} --- planned for future work.
\end{itemize}

\textbf{Metrics.}
Cumulative search time (ms), insert time (ms), delete time (ms),
maintenance time (ms), total time (ms), mean recall, and final partition count.

% ─────────────────────────────────────────────────────────────────────────────
\subsection{Refinement Dominates Maintenance Cost}
\label{sec:eval_ablation}
% Answers: why does selective refinement matter at all?
% ─────────────────────────────────────────────────────────────────────────────

% [SINGLE-RUN] Table from seed 9299.
\begin{table}[h]
\centering
\caption{Ablation: disabling refinement cuts maintenance 14$\times$
  but drops recall 3 points. \textit{Single-run (seed 9299).}}
\label{tab:ablation}
\begin{tabular}{lrrr}
  \toprule
  Configuration & Maintain & Total & Recall \\ \midrule
  Quake              & 103,926 & 374,684 & 0.909 \\
  Quake-NoCost       &  62,816 & 386,286 & 0.900 \\
  Quake-NoRefine     &   7,263 & 289,375 & 0.879 \\
  Quake-NoReject     &  99,246 & 348,580 & 0.768 \\
  \bottomrule
\end{tabular}
\end{table}

Disabling refinement (NoRefine) cuts maintenance from 103,926 to
7,263\,ms ($14\times$) but drops recall from 0.909 to 0.879.
Disabling the cost-based rejection mechanism (NoReject) collapses recall
to 0.768. These results establish refinement as the dominant maintenance
cost and motivate its selective application: eliminating it blindly
produces unacceptable recall loss.

% ─────────────────────────────────────────────────────────────────────────────
\subsection{Score-Based Refinement Reduces Maintenance by 38--45\%
            While Preserving Recall}
\label{sec:eval_main}
% Answers RQ1: does FlexiRefine reduce maintenance while preserving recall?
% ─────────────────────────────────────────────────────────────────────────────

\textbf{Policy exploration (single-run, seed 9299).}
Table~\ref{tab:density} shows results for a sweep of score-based
policies with density-normalized staleness
$D(p) = m_p / |p|$ at varying budget $K$.

% [SINGLE-RUN]
\begin{table}[h]
\centering
\caption{Score-based refinement sweep with density-normalized staleness.
  All configs use split threshold 240\,ns, radius=50, iterations=2.
  \textit{Single-run (seed 9299).}}
\label{tab:density}
\begin{tabular}{lrrrr}
  \toprule
  Config & Maintain & Total & Recall & Partitions \\
  \midrule
  Quake                   & 72,176 & 370,609 & 0.9093 & 2319 \\
  Score50-RawMutations    & 35,292 & 381,188 & 0.9018 & 1333 \\
  Density30-NormMutations & 26,906 & 381,575 & 0.8960 & 1290 \\
  Density50-NormMutations & 34,726 & 375,284 & 0.8999 & 1333 \\
  Density70-NormMutations & 41,091 & 377,188 & 0.9041 & 1364 \\
  Density100-NormMutations& 41,596 & 379,794 & 0.9041 & 1335 \\
  \bottomrule
\end{tabular}
\end{table}

Density70 achieves recall 0.9041 with maintenance $-43\%$ versus full
Quake and total time within 1.8\%. Density50 hits the recall target
(0.8999) with the lowest maintenance of above-target configs.
Both configurations improve over raw mutation scoring at the same $K$.

\textbf{Exponent structure (single-run).}
Table~\ref{tab:grid} explores the general score family
$S_{\beta,\gamma}(p) = (h_p+1)^\beta D(p)^\gamma$
over a $3\times 3$ grid of exponent values at $K=70$ with
density normalization enabled.

% [SINGLE-RUN]
\begin{table}[h]
\centering
\caption{Beta/gamma exponent grid for Density70 ($K=70$,
  density normalization). \textit{Single-run (seed 9299).}}
\label{tab:grid}
\begin{tabular}{ccrrr}
  \toprule
  $\beta$ & $\gamma$ & Maintain & Total & Recall \\ \midrule
  0.5 & 0.5 & 37,805 & 376,547 & 0.9033 \\
  0.5 & 1.0 & 45,214 & 393,305 & 0.9029 \\
  0.5 & 1.5 & 54,977 & 382,108 & 0.9049 \\
  \textbf{1.0} & \textbf{0.5} & \textbf{37,382} & \textbf{370,039} & 0.9035 \\
  1.0 & 1.0 & 50,030 & 393,363 & \textbf{0.9052} \\
  1.0 & 1.5 & 39,504 & 379,727 & 0.9018 \\
  1.5 & 0.5 & 40,455 & 379,769 & 0.9016 \\
  1.5 & 1.0 & 35,535 & 374,353 & 0.9013 \\
  1.5 & 1.5 & 37,364 & 377,434 & 0.9018 \\
  \bottomrule
\end{tabular}
\end{table}

Two column-level patterns are robust:
(1)~The $\gamma=0.5$ column consistently produces lower total runtime
(370k--380k) than $\gamma=1.0$ (374k--393k), suggesting sublinear
staleness compression reduces overhead.
(2)~The $\beta=1.5$ row has consistently lower recall (0.9013--0.9018),
confirming that super-linear hit amplification is detrimental.
Individual cell comparisons are within single-run noise and are not
treated as final; see~\S\ref{sec:eval_seeds} for multiseed evidence.

% ─────────────────────────────────────────────────────────────────────────────
\subsection{Coverage-Only and Naive Policies Fail to Preserve Recall}
\label{sec:eval_negative}
% Answers RQ2: are score-based policies better than naive alternatives?
% ─────────────────────────────────────────────────────────────────────────────

\textbf{Split-threshold tuning.}
Higher split thresholds reduce maintenance aggressively but recall
cannot reach 0.90 regardless of threshold (Table~\ref{tab:split}).
Reducing split frequency alone is insufficient; selective post-split
repair is also required.

% [SINGLE-RUN]
\begin{table}[h]
\centering
\caption{Split-threshold sweep. Maintenance falls but recall stays
  below target. \textit{Single-run.}}
\label{tab:split}
\begin{tabular}{lrrr}
  \toprule
  Threshold (ns) & Maintain & Total & Recall \\ \midrule
  150 &  9,165 & 323,147 & 0.8826 \\
  200 &  5,861 & 334,535 & 0.8845 \\
  240 &  5,222 & 344,766 & 0.8860 \\
  300 &  8,774 & 366,886 & 0.8837 \\
  \bottomrule
\end{tabular}
\end{table}

\textbf{Global budget cap without value ranking.}
Applying a global top-$K$ candidate cap without a ranking signal
refines fewer partitions, not better ones. All cap values keep recall
below 0.90 (Table~\ref{tab:topk}).

% [SINGLE-RUN]
\begin{table}[h]
\centering
\caption{Global top-K candidate cap without score ranking
  (radius=50, iterations=2). \textit{Single-run.}}
\label{tab:topk}
\begin{tabular}{lrrr}
  \toprule
  Policy & Maintain & Total & Recall \\ \midrule
  Refine-all &  39,626 & 362,887 & 0.8991 \\
  Top-5      &  22,189 & 393,564 & 0.8872 \\
  Top-10     &  22,181 & 393,253 & 0.8925 \\
  Top-20     &  29,363 & 384,396 & 0.8965 \\
  \bottomrule
\end{tabular}
\end{table}

\textbf{Per-split coverage budgeting.}
\superseded{Note: earlier numbers (PerSplit3-Top10 maintain 21,584,
recall 0.8903) were produced by a code bug in which
\texttt{refinement\_candidates\_per\_split} was not applied.
Corrected numbers below use commit \texttt{ea4de23}.}

Per-split budgeting allocates $b$ candidate slots per split centroid
before the global union.
After the bug fix, no per-split configuration reaches recall 0.90
(Table~\ref{tab:persplit}).
The best result is 0.8945, well below global unranked refinement (0.9039).

% [SINGLE-RUN, corrected]
\begin{table}[h]
\centering
\caption{Per-split budgeting after correctness fix.
  No configuration reaches recall 0.90. \textit{Single-run (seed 9299).}}
\label{tab:persplit}
\begin{tabular}{lrrr}
  \toprule
  Configuration & Maintain & Total & Recall \\ \midrule
  Quake                   & 72,701 & 365,861 & 0.9068 \\
  NoRefine-Split240       &  9,422 & 388,623 & 0.8855 \\
  RefineAll (radius=50)   & 42,824 & 378,127 & 0.9039 \\
  PerSplit3-Top20         & 29,996 & 393,098 & 0.8894 \\
  PerSplit5-Top20         & 26,616 & 397,353 & 0.8945 \\
  PerSplit5-Top30         & 28,413 & 392,926 & 0.8945 \\
  \bottomrule
\end{tabular}
\end{table}

Per-split forcing breaks global value ranking. When split centroids
are spatially close, their neighborhoods overlap; the per-centroid cap
blocks globally high-value candidates that do not rank first from every
nearby centroid. Coverage guarantees do not substitute for a strong
importance--staleness signal.

\textbf{Summary of negative results.}
Size-only filtering prioritizes large stable partitions over small stale
ones. Distance-only filtering identifies structurally affected candidates
but not queried ones. Budget caps without ranking refine fewer
partitions, not better partitions. Per-split coverage enforces spatial
coverage but breaks global value ordering. Hit-only scoring captures
workload importance but ignores structural degradation. Super-linear
hit amplification ($\beta > 1$) over-concentrates budget on a few hot
partitions. All of these failure modes are avoided by jointly scoring
with workload importance and structural staleness.

% ─────────────────────────────────────────────────────────────────────────────
\subsection{Static IVF Avoids Maintenance but Pays 4\texttimes{} Higher
            Total Runtime}
\label{sec:eval_baselines}
% Answers RQ3: is dynamic maintenance necessary?
% ─────────────────────────────────────────────────────────────────────────────

Table~\ref{tab:baseline_comparison} compares \system configurations
against Faiss-IVF under the full 3-seed dynamic workload.
Numbers are mean\,$\pm$\,std across seeds $\{9299, 42, 12345\}$.

% [CONFIRMED 3-seed]
\begin{table}[h]
\centering
\caption{Baseline comparison: \system vs.\ Faiss-IVF (no maintenance).
  Mean\,$\pm$\,std across three seeds. \textbf{[CONFIRMED]}}
\label{tab:baseline_comparison}
\begin{tabular}{lrrr}
  \toprule
  Config & Maintain (ms) & Total (ms) & Recall \\ \midrule
  Quake
    & $70{,}682 \pm 15{,}431$
    & $309{,}741 \pm 66{,}555$
    & $0.9118 \pm 0.0134$ \\
  Score50-Raw-Linear
    & $38{,}579 \pm 7{,}528$
    & $321{,}162 \pm 64{,}340$
    & $0.9017 \pm 0.0192$ \\
  Density70-Linear
    & $43{,}735 \pm 11{,}913$
    & $323{,}591 \pm 54{,}505$
    & $0.9039 \pm 0.0165$ \\
  D70-B0.5-G0.5
    & $42{,}224 \pm 4{,}609$
    & $316{,}609 \pm 54{,}402$
    & $0.9064 \pm 0.0184$ \\
  Faiss-IVF (nprobe=20)
    & $0 \pm 0$
    & $1{,}315{,}646 \pm 782{,}124$
    & $0.9269 \pm 0.0668$ \\
  \bottomrule
\end{tabular}
\end{table}

Faiss-IVF spends nothing on maintenance but incurs approximately
$4.25\times$ higher total runtime than Quake and $4.16\times$ higher
than D70-B0.5-G0.5. Its insert and delete latencies are low (9,031 and
732\,ms respectively, versus ~18,000 and ~23,000 for Quake) because no
centroid updates occur, but this advantage is more than offset by search
cost: as the skewed workload inserts vectors into outdated centroid
neighborhoods, some partitions grow much larger than others, and
scanning a fixed nprobe=20 becomes progressively more expensive.
Critically, recall variance is $\pm$0.0668 --- one seed produced recall
0.8553, below the 0.90 target, showing that a static index cannot
reliably maintain recall under a dynamic workload.

\textbf{Faiss-IVF framing.}
Faiss-IVF should be read as a \emph{degradation baseline}: it quantifies
the cost of forgoing maintenance entirely. The 70,682\,ms maintenance
savings Quake spends on refinement costs over 1,000,000\,ms in added
search time over the workload --- a $14\times$ exchange ratio. Dynamic
maintenance is necessary.

\textbf{LIRE-style baseline.}
\todo{Replace with confirmed 3-seed results from
  \texttt{sift1m\_split240\_lire\_baseline\_3seed}.
  Expected result: LIRE preserves recall between
  Faiss-IVF and Quake, but falls below 0.90 without K-means refinement,
  confirming that reassignment alone is insufficient.
  Key comparison: Quake-NoRefine (fixed splits) vs.\
  LIRE (size-based splits, no K-means) vs.\
  FlexiRefine (selective K-means refinement).}

% ─────────────────────────────────────────────────────────────────────────────
\subsection{Score-Based Refinement Generalizes Across Workload Seeds}
\label{sec:eval_seeds}
% Answers RQ4: are results robust across seeds?
% ─────────────────────────────────────────────────────────────────────────────

Table~\ref{tab:confirmation} reports the primary \system result:
mean\,$\pm$\,std across seeds $\{9299, 42, 12345\}$ for the six
strongest configurations from the exploratory sweeps.

% [CONFIRMED 3-seed, PRIMARY TABLE]
\begin{table}[h]
\centering
\caption{3-seed confirmation of \system selective-refinement policies.
  All score-based configs preserve mean recall $\geq$ 0.90.
  \textbf{[CONFIRMED]}}
\label{tab:confirmation}
\begin{tabular}{lrrr}
  \toprule
  Config & Maintain (ms) & Total (ms) & Recall \\ \midrule
  Quake
    & $67{,}658 \pm 6{,}505$
    & $308{,}584 \pm 54{,}185$
    & $0.9084 \pm 0.0149$ \\
  Score50-Raw-Linear
    & $36{,}620 \pm 8{,}825$
    & $320{,}085 \pm 56{,}744$
    & $0.9026 \pm 0.0181$ \\
  Density70-Linear
    & $42{,}946 \pm 5{,}238$
    & $317{,}650 \pm 55{,}656$
    & $0.9033 \pm 0.0183$ \\
  D70-B0.5-G1.0
    & $46{,}038 \pm 15{,}398$
    & $332{,}935 \pm 79{,}071$
    & $0.9057 \pm 0.0174$ \\
  D70-B1.0-G0.5
    & $43{,}441 \pm 5{,}644$
    & $321{,}090 \pm 50{,}160$
    & $0.9030 \pm 0.0206$ \\
  D70-B0.5-G0.5
    & $40{,}941 \pm 8{,}024$
    & $316{,}805 \pm 53{,}578$
    & $0.9024 \pm 0.0215$ \\
  \bottomrule
\end{tabular}
\end{table}

\textbf{Confirmed findings.}
\begin{itemize}[leftmargin=*]
  \item All score-based selective-refinement configs preserve mean recall
    $\geq$ 0.90. Recall means range from 0.9024 to 0.9057.
  \item Maintenance reductions of 32--46\% versus Quake are confirmed.
    Score50-Raw-Linear: $-45.9\%$; D70-B0.5-G0.5: $-39.5\%$;
    Density70-Linear: $-36.5\%$.
  \item Total runtime overhead stays within 4\% of Quake for all configs
    except D70-B0.5-G1.0 ($+7.9\%$).
  \item Density70-Linear has the lowest maintenance variance (CV=12.2\%),
    making it the most behaviorally predictable configuration.
\end{itemize}

\textbf{Policy-class claim versus hyperparameter claim.}
The exact ranking within the score-based policy class is not stable at
$n=3$ seeds. The maintenance difference between Density70-Linear
(42,946\,ms) and D70-B1.0-G0.5 (43,441\,ms) is 495\,ms --- smaller than
either config's standard deviation. This motivates claims at the
\emph{policy-class level}: score-based refinement with a joint
importance--staleness signal robustly outperforms all naive baselines,
independent of the specific score parameters chosen.

% ─────────────────────────────────────────────────────────────────────────────
\subsection{Selecting an Operating Point}
\label{sec:eval_operating_point}
% Answers RQ5: which configuration should a user choose?
% ─────────────────────────────────────────────────────────────────────────────

Different deployments have different priorities. Based on the confirmed
3-seed results, we recommend:

\begin{itemize}[leftmargin=*]
  \item \textbf{Recall-critical users.}
    Use \textbf{Quake} full refinement.
    Mean recall 0.9084, lowest variance ($\pm$0.0134).
    Accept full maintenance overhead (70,682\,ms mean) to guarantee
    the strongest recall.

  \item \textbf{Maintenance-constrained users.}
    Use \textbf{Score50-Raw-Linear}: $S(p) = (h_p+1)\,m_p$, $K=50$.
    Mean maintenance 36,620\,ms ($-45.9\%$ versus Quake), recall 0.9026.
    Warning: higher maintenance variance (CV=24.1\%); behavior is less
    predictable than density-normalized alternatives.

  \item \textbf{Balanced users.}
    Use \textbf{D70-B0.5-G0.5}: $S(p) = (h_p+1)^{0.5}(m_p/|p|)^{0.5}$,
    $K=70$.
    Best single configuration on total runtime ($+2.2\%$ vs Quake),
    maintenance $-39.5\%$, recall 0.9064.
    Lowest maintenance standard deviation (CV=10.9\%).

  \item \textbf{Simplicity and stability.}
    Use \textbf{Density70-Linear}: $S(p) = (h_p+1)(m_p/|p|)$, $K=70$.
    Simplest density formulation, lowest maintenance CV (12.2\%),
    competitive recall (0.9033), total overhead $+2.9\%$.
    Recommended default for deployments that prioritize
    predictable behavior over absolute minimum maintenance.
\end{itemize}

\textbf{What to avoid.}
Do not use a static IVF index (Faiss-IVF) for dynamic workloads:
it pays $4.25\times$ higher total runtime and cannot reliably maintain
the recall target across different workload seeds.
Do not use per-split coverage budgeting without a ranking signal:
it reduces maintenance but fails to preserve recall.

% ─────────────────────────────────────────────────────────────────────────────
\subsection{Comparison with Graph-Based and Optimized Baselines}
\label{sec:eval_graph}
% Answers RQ5 extension: optional comparison with HNSW, ScaNN
% ─────────────────────────────────────────────────────────────────────────────

\todo{Add HNSW comparison on insert-only or delete-free workload variant.
  FaissHNSW wrapper exists (\texttt{src/python/index\_wrappers/faiss\_hnsw.py})
  but requires a workload without delete operations since the Faiss HNSW
  implementation does not support vector removal.
  Key framing: HNSW achieves strong search quality in static or
  insert-only settings but cannot handle the full dynamic workload.
  Suggest running HNSW on a 70\% query / 30\% insert variant.}

\todo{Add ScaNN comparison. Wrapper exists (\texttt{src/python/index\_wrappers/scann.py}).
  Requires \texttt{pip install scann} on CloudLab.
  ScaNN supports inserts and deletes; can run on the full dynamic workload.}

% ============================================================
\section{Discussion}
\label{sec:discussion}
% ============================================================

\subsection{Run Variance and K-Means Non-Determinism}

Refinement quality depends on K-means centroid initialization, which is
stochastic and not controlled by the workload seed. Over a 1000-operation
dynamic workload, small early-cycle differences compound into meaningful
behavioral divergence. We observe maintenance cost coefficients of
variation between 10\% and 33\% across three seeds, and recall standard
deviations of $\pm$0.015--0.022. This motivates reporting
mean\,$\pm$\,std rather than single-run tables, and framing primary
claims at the policy-class level rather than asserting a specific
$(\beta, \gamma)$ pair as optimal.

\subsection{Policy-Class Claims versus Hyperparameter Claims}

The three-seed data establishes that score-based selective refinement
robustly preserves mean recall above 0.90 while reducing maintenance
by 32--46\%. Differences between specific configurations within the
score-based class (e.g., Density70-Linear versus D70-B1.0-G0.5) are
within run-to-run variance at $n=3$ seeds. Claims about optimal
$\beta$, $\gamma$, or exact normalization require $n \geq 5$ seeds
or deterministic K-means initialization.

\subsection{Scope of Current Results}

All confirmed results use SIFT1M with a 30/20/50 insert/delete/query
workload mix. The policy framework is architecture-agnostic, but
empirical validation on larger datasets (SIFT10M, MSTuring) and
different workload mixes is needed before general claims about
dataset sensitivity can be made.

\subsection{Baseline Coverage}

Faiss-IVF has been evaluated and conclusively establishes that dynamic
maintenance is necessary: a static index pays $4.25\times$ higher total
runtime and fails the recall target on one of three seeds.
LIRE-style maintenance (size-based splits, no K-means refinement) is
in progress --- preliminary results are expected to show that dynamic
splits alone are insufficient without centroid-quality repair.
Graph-based baselines (HNSW, ScaNN) are planned and will address whether
the recall--total-runtime tradeoff achieved by \system is competitive
with non-IVF approaches.

% ============================================================
\section{Conclusion}
\label{sec:conclusion}
% ============================================================

This paper argues that dynamic vector indexing should be treated as a
controllable system rather than a fixed algorithm. Refinement is a
powerful control knob: it strongly affects both maintenance cost and
search quality, and its selective application exposes a useful
tradeoff frontier.

\system frames refinement as a resource allocation problem. The key
insight, confirmed across three independent workloads, is that
refinement should target partitions that are both important to queries
and structurally stale. Naive policies that filter by size, distance,
or local coverage fail to preserve the recall target while reducing
maintenance. Score-based selective refinement using a joint
importance--staleness signal reduces maintenance by \textbf{36--46\%}
while preserving \textbf{mean recall above 0.90} with total runtime
within \textbf{4\% of the full-refinement baseline}.

Among tested configurations, density-normalized staleness
($D(p) = m_p / |p|$) provides the most behaviorally stable operating
point (maintenance CV=12.2\%), making it the recommended default.
Raw mutation scoring achieves lower absolute maintenance but with
higher variance. The precise optimal exponents $(\beta, \gamma)$ in
the general score family require further confirmation.

% ============================================================
\section{Future Work}
\label{sec:future}
% ============================================================

\begin{enumerate}[leftmargin=*]
  \item \textbf{External baselines.}
    Add Faiss-IVF (minimum required), a LIRE/SPFresh-style
    baseline for dynamic-maintenance comparison, and one graph-based
    baseline (HNSW or DiskANN). Wrapper implementations exist.

  \item \textbf{More seeds.}
    Extend to $n=5$ seeds (add 1337 and 2024) to tighten standard
    deviations and enable more defensible hyperparameter comparisons.

  \item \textbf{Deterministic K-means seed control.}
    Exposing a random seed parameter in the C++ refinement K-means
    would eliminate run variance entirely and simplify reproducibility.
    Requires a small change to \texttt{refine\_partitions()} in
    \texttt{partition\_manager.cpp}.

  \item \textbf{Larger workloads.}
    Evaluate on SIFT10M and MSTuring-10M. The system supports both
    datasets; experiment configs would be straightforward additions.

  \item \textbf{Gamma fine sweep.}
    With $\beta$ fixed at 1.0, sweep $\gamma \in \{0.25, 0.5, 0.75,
    1.0\}$ to locate the optimal staleness exponent more precisely.

  \item \textbf{Adaptive budget control.}
    Use observed recall drift to adjust the top-$K$ budget
    dynamically, targeting a recall threshold rather than a fixed count.
    One step toward the adaptive-budget future work outlined in
    the paper.
\end{enumerate}

% ============================================================
\bibliographystyle{ACM-Reference-Format}
\bibliography{references}
% ============================================================

\end{document}
